\subsubsection{Synchronization Time Under Random Real Input: Subset--Memory Absorbing Markov Chain (Closed Form)}\label{app:real-input-40-sync-time}
This section promotes the ``existence of a shortest reset word $0^5$'' from Proposition \ref{prop:real-input-40-reset-word} to a random performance metric: under the Parry maximum entropy measure of the real input $X\times X$, how quickly is the internal synchronization state forced to become unique when observing the input stream $(x_k,y_k)$?

\begin{definition}[Subset--memory absorbing chain and synchronization time]\label{def:real-input-40-sync-time-absorbing}
View the synchronization kernel as a deterministic automaton \(\delta:Q\times\{0,1,2\}\to Q\) (ignoring the output). Let the input be the digit-wise sum \(d_k=x_k+y_k\) of two golden mean orbits \(x,y\in X\), and denote the input memory as \(m_k:=(x_{k-1},y_{k-1})\in\{00,01,10,11\}\).
\par
If the initial synchronization state is unknown but only known to lie in some set \(S_0\subseteq Q\), then after reading an input block of length \(n\), the ``possible synchronization state set'' is deterministically updated via the subset construction to
\[
S_n=\delta(S_0,d_1d_2\cdots d_n)\subseteq Q.
\]
The minimum number of steps to reach \(|S_n|=1\),
\[
\tau:=\min\{n\ge 1:\ |S_n|=1\},
\]
is called the \textbf{synchronization time} (absorption time).
Under the Parry measure, \((S_n,m_n)\) forms a finite-state absorbing Markov chain: \(m_n\) is the 4-state skeleton chain of \(X\times X\), and \(S_{n+1}\) is given by the deterministic map of \(d_{n+1}\) applied to \(S_n\).
\end{definition}

\begin{proposition}[Mean/variance of synchronization time, exponential tail rate, and quantiles]\label{prop:real-input-40-sync-time-moments}
Under the absorbing chain of Definition \ref{def:real-input-40-sync-time-absorbing}, consider two natural initial uncertainty scenarios:
\begin{itemize}
  \item \textbf{worst}: the initial synchronization state is only known to be arbitrary in \(Q\), so \(S_0=Q\);
  \item \textbf{essential}: the initial $(s,m)$ is only known to lie in the essential $20$-state core of Proposition \ref{prop:real-input-40-essential-20}, so \(S_0\) is taken as the corresponding set of allowed synchronization states for each memory class \(m\) ($6/5/5/4$ stratification).
\end{itemize}
Take \(m_0\) as the Parry stationary marginal distribution of \(X\times X\) (Proposition \ref{prop:real-input-40-input-memory-marginal}).
Then \(\mathbb{E}[\tau]\) and \(\mathrm{Var}(\tau)\) are given in \(\mathbb{Q}(\sqrt5)\), and the absorption tail satisfies exponential decay with transient spectral radius
\[
\rho=\varphi^{-1}=\frac{\sqrt5-1}{2}.
\]
Specific numerical values and quantile summaries are provided in Table \ref{tab:real-input-40-sync-time-absorbing-chain} (automatically generated by the script, including exact closed forms and reproducible quantification).
\end{proposition}

\input{sections/generated/tab_real_input_40_sync_time_absorbing_chain}

\begin{remark}[Reproducible output and computational caliber]\label{rem:real-input-40-sync-time-repro}
Table \ref{tab:real-input-40-sync-time-absorbing-chain} is generated by the script \path{scripts/exp_real_input_40_sync_time_absorbing_chain.py}.
Therein, \(\mathbb{E}[\tau]\) and \(\mathrm{Var}(\tau)\) are obtained by exactly solving the absorbing chain linear equations
\(
(I-Q)t=\mathbf{1}
\)
and
\(
(I-Q)u=\mathbf{1}+2Qt
\)
over \(\mathbb{Q}(\sqrt5)\) (\(Q\) denotes the transient submatrix), while the quantiles are obtained by numerical recursion \(v_{n+1}=v_n Q\) on the transient chain.
\end{remark}
